\documentclass[11pt,a4paper]{article}

\usepackage[utf8]{inputenc}
\usepackage[T2A]{fontenc}
\usepackage[russian,english]{babel}
\usepackage{lmodern}
\usepackage[margin=2.2cm]{geometry}
\usepackage{booktabs}
\usepackage{longtable}
\usepackage{array}
\usepackage{hyperref}
\usepackage{enumitem}
\usepackage{parskip}
\usepackage{titlesec}
\usepackage{xcolor}
\usepackage{listings}

\hypersetup{
    colorlinks=false,
    hidelinks
}

\titleformat{\section}{\large\bfseries}{\thesection}{0.6em}{}
\titleformat{\subsection}{\normalsize\bfseries}{\thesubsection}{0.6em}{}

\setlist[enumerate,1]{label=\Roman*., leftmargin=2.2em}
\setlist[enumerate,2]{label=\alph*), leftmargin=2.0em}

\lstset{
    basicstyle=\ttfamily\small,
    breaklines=true,
    frame=none,
    columns=fullflexible,
    keepspaces=true,
    showstringspaces=false
}

\begin{document}

\begin{center}
    {\LARGE\bfseries NGSolve Agent - журнал разработки}
\end{center}

\vspace{1em}

Этот файл - рабочий журнал проекта: что сделано, почему, как проверено и что делать дальше. Записи идут от новых к старым; после каждого существенного изменения кода или запуска добавляется новая запись.

\section{Запись 2026-11-10 - реальные LLM-адаптеры для D1/D4}

\subsection{Что изменено}

\begin{enumerate}
    \item Добавлен \texttt{src/ngsolve\_agent/llm\_adapter.py}: один JSON-only адаптер поддерживает Ollama local chat и OpenAI-compatible chat completions. Модель задаётся через переменные окружения; API key не попадает в отчёт и исходники.
    \item Добавлен \texttt{src/ngsolve\_agent/agent\_cli.py}: запускает D1 или D4 на обычном пользовательском запросе и пишет запрос, архитектуру, provider/model, trace, время/usage LLM-вызовов, solver status и путь к численному отчёту.
    \item D4 теперь сохраняет шаги interpret, validate, solve, evaluate в trace.
    \item В README добавлена конфигурация запуска локального Ollama и OpenAI-compatible API.
\end{enumerate}

\subsection{Зачем и границы}

Раньше D1/D4 принимали только тестовые callback-заглушки. Теперь у проекта есть реальный API-клиентский путь, но FEM остаётся детерминированным и валидируемым: модель может только преобразовать запрос в \texttt{PoissonTask} либо выбрать один из allowlisted D1-инструментов. Произвольный Python-код не генерируется и не запускается. Этот первый LLM-пилот пока ограничен поддерживаемой Poisson-задачей; перенос на Stokes/elasticity/Maxwell - следующий инженерный этап.

\subsection{Проверка и ограничение среды}

Ruff прошёл для \texttt{examples}, \texttt{src} и \texttt{tests}; \texttt{compileall} завершился без ошибок. Фактический LLM-вызов пока не запускался: локальный Ollama executable/endpoint не обнаружен, \texttt{NGSOLVE\_LLM\_API\_KEY} не задан. Значения секретов не читались и не выводились. После настройки доступного endpoint запускать через:

\begin{lstlisting}
uv run python -m ngsolve_agent.agent_cli --architecture d4 --request "..."
uv run python -m ngsolve_agent.agent_cli --architecture d1 --request "..."
\end{lstlisting}

Сравнительный LLM-run нельзя считать выполненным до подключения модели; в отчёте это будет видно по \texttt{llm\_calls} и полю \texttt{status}.

\subsection{Состояние локальной установки (проверено повторно)}

\begin{enumerate}
    \item В Windows-среде Codex \texttt{ollama} CLI не найден, а \texttt{http://127.0.0.1:11434/api/tags} не отвечает. Codex не смог установить Ollama: \texttt{winget} завис на загрузке; прямое обращение из терминала к GitHub не прошло из-за недоступности DNS в управляемом окружении. Пользовательское Windows-окружение может иметь другую сетевую политику.
    \item Проверка RAM/GPU через CIM и \texttt{systeminfo} также недоступна из управляемого процесса. Поэтому выбран небольшой Qwen3.5 4B для первого прогона без заявления о скорости; официальная карточка Ollama указывает размер около 3.3-4.0 GB.
    \item Для запуска из обычного Windows Terminal пользователь должен установить Ollama официальной командой Windows, затем выполнить \texttt{ollama run qwen3.5:4b} и запускающую команду D4 из README. Имя службы \texttt{OllamaService} не проверено и не требуется для штатного запуска из приложения Ollama.
    \item В предоставленном внешнем ответе \texttt{d4} ошибочно прочитано как степень/уровень метода; в проекте D4 - архитектурная ветка с фиксированным workflow LangGraph.
    \item Пока Ollama endpoint отсутствует, выполненных LLM-вызовов и сравнительных результатов нет. Не подменять это заглушками в экспериментальном отчёте.
\end{enumerate}

\subsection{Подготовка controlled сравнения}

\begin{enumerate}
    \item Добавлен фиксированный набор \texttt{data/agent\_poisson\_cases.json} (P1-P5): пять естественно-языковых запросов с эталонными значениями параметров \texttt{PoissonTask}.
    \item Добавлен \texttt{examples/evaluate\_agent\_architectures.py}: запускает выбранный набор на одинаковой модели через D1 и D4, проверяет точное извлечение параметров отдельно от численного \texttt{passed}, считает вызовы/время LLM и пишет comparison JSON.
    \item Первый запуск \texttt{--case P1} предназначен как дешёвая проверка интеграции, прежде чем запускать все десять сочетаний архитектуры и задания.
    \item После добавления comparison runner Ruff прошёл по \texttt{examples}, \texttt{src}, \texttt{tests}, \texttt{compileall} прошёл; \texttt{--help} обеих новых CLI-команд отработал.
    \item Проверка окружения: \texttt{ollama} не найден, \texttt{http://127.0.0.1:11434/api/tags} не отвечает, API-key переменная не настроена. Поэтому реальных model calls пока нет.
\end{enumerate}

\section{Текущее состояние}

\begin{enumerate}
    \item Проект: Windows, Python 3.12, NGSolve 6.2.2608; зависимости закреплены в \texttt{pyproject.toml} и \texttt{uv.lock}.
    \item Запускаются четыре детерминированных FEM benchmark-а: Poisson, нелинейная Neo-Hookean балка, Stokes flow вокруг цилиндра и Maxwell curl-curl eigenproblem на Fichera-геометрии.
    \item \texttt{examples/run\_benchmark\_suite.py} выполняет все четыре задачи и записывает единый отчёт с проверками в \texttt{results/benchmark\_runs/}.
    \item D1 и D4 реализованы как каркасы с внедряемыми callback-ами. Реальная LLM-модель ещё не подключена; MCP-интерфейс ещё не добавляли.
    \item Дообучение модели не планируется на текущем этапе. Локальная модель - возможный будущий вариант сравнения, а не требование Александра.
\end{enumerate}

\section{Запись 2026-11-10 - единый запуск четырёх benchmark-ов}

\subsection{Что изменено}

\begin{enumerate}
    \item Добавлен \texttt{examples/run\_benchmark\_suite.py}. Он запускает B0 Poisson, нелинейную упругость, Stokes и Maxwell последовательно, выполняет задачно-зависимые проверки и пишет агрегированный \texttt{suite.json} вместе с per-case JSON и графиками.
    \item Добавлена зависимость SciPy в \texttt{pyproject.toml} и \texttt{uv.lock}: NGSolve PINVIT импортирует \texttt{scipy.linalg} при решении Maxwell eigenproblem.
    \item В README добавлена команда запуска всего набора.
\end{enumerate}

\subsection{Почему}

До этого задачи запускались разными командами и имели разные форматы результата. Общий запуск нужен, чтобы позже сравнивать B0, D1 и D4 на тех же задачах и сопоставлять проверки, ошибки и время выполнения.

\subsection{Запуск и результат}

Команда:

\begin{lstlisting}
.venv\Scripts\python.exe examples\run_benchmark_suite.py
\end{lstlisting}

Итог: все 4 из 4 задач прошли текущие проверки (\texttt{all\_passed: true}), NGSolve 6.2.2608, Python 3.12.16. Отчёт этого запуска: \texttt{results/benchmark\_runs/suite-20261110T130937820324Z/suite.json}.

Ключевые результаты:

\begin{enumerate}
    \item Poisson: L2 error \texttt{1.6335e-3}; свободная невязка \texttt{1.1871e-15}.
    \item Neo-Hookean cantilever: все 50 шагов нагрузки сошлись; свободная невязка \texttt{1.1568e-13}.
    \item Stokes: поток на входе и выходе \texttt{0.41}; относительный дисбаланс \texttt{0}; свободная невязка \texttt{1.0349e-14}. L2-норма дивергенции \texttt{0.3157}; это измеряемая величина, не нулевой pass/fail порог для Taylor-Hood.
    \item Maxwell/Fichera: 12 положительных собственных значений, 42,103 степени свободы; первая собственная величина \texttt{3.2202713671}, относительное отличие от значения в tutorial \texttt{3.88e-6}.
\end{enumerate}

Ruff прошёл для \texttt{examples}, \texttt{src} и \texttt{tests}. Полный suite был реально исполнен; тесты \texttt{pytest} в этой записи повторно не запускались.

\subsection{Технические решения и замечания}

\begin{enumerate}
    \item Первоначальный упрощённый 2D Maxwell-вариант не дал ожидаемую первую моду и был отброшен. В журналируемый набор вошёл опубликованный NGSolve 3D Fichera-пример, чтобы сравнивать с его устойчивым справочным спектром.
    \item Для установки SciPy использовалась временная папка \texttt{.uv-cache} в рабочей папке; после установки её удалили. OneDrive подтвердил удаление 334 файлов этого кеша. Это были скачиваемые временные данные, не исходники проекта и не результаты.
    \item Smart App Control ранее блокировал неподписанные исполняемые файлы окружения; пользователь отключил блокировку самостоятельно. Код проекта обходы системной защиты не добавляет.
\end{enumerate}

\section{Сделано ранее}

\subsection{Исследовательский замысел и ограничения}

\begin{enumerate}
    \item Прочитан \texttt{NGSolve\_Agent\_Research\_Blueprint.pdf}. Основной вопрос проекта - надёжность и воспроизводимость архитектур управления научным расчётом.
    \item Зафиксированы сравниваемые архитектуры: B0 - ручной детерминированный решатель; D1 - модель выбирает типизированные инструменты; D4 - фиксированный workflow LangGraph с проверками и ограниченным восстановлением. D3/RAG позже, D2 опционально.
    \item По словам Александра Хейнленки нет требования обязательно использовать локальную модель или обучать новую. Начинаем с агента и реальных расчётов; провайдер и local-vs-API сравнение выбираются отдельным этапом.
\end{enumerate}

\subsection{Кодовая основа}

\begin{enumerate}
    \item Созданы Poisson B0-решатель, Pydantic-схемы, независимая численная проверка, JSON-отчёт и изображение решения.
    \item Созданы D1 и D4 как provider-neutral каркасы с callback-ами; их модельные решения пока проверялись тестовыми заглушками, не подключённой LLM.
    \item Исправлены вызовы NGSolve API для вычисления H1-ошибки и свободной невязки.
    \item \texttt{uv.lock} и \texttt{.venv} используются для воспроизводимого Windows-окружения.
\end{enumerate}

\section{Следующие шаги}

\begin{enumerate}
    \item Добавить для Stokes ещё один уровень сетки и оценивать изменение divergence/давления, потому что сейчас есть только два уровня.
    \item Привести проверки и метрики четырёх задач к формально описанному общему контракту (task ID, параметры, метрики, evaluator version, артефакты), не стирая различия между физическими задачами.
    \item Добавить реальные задания NGSolve из учебных материалов: уточнить набор задач, эталонные/проверочные величины и допустимые варианты формулировки.
    \item Подключить один реальный LLM provider к D1 и D4, сохраняя текущий детерминированный FEM solver и evaluator; секреты не хранить в репозитории.
    \item Сравнить B0/D1/D4 на одном наборе задач и ограниченном числе повторов. Сохранять model/version, входное условие, выбранные инструменты, trace, повторы, ошибки, численные проверки и время.
    \item Только после первых сравнений решить, нужен ли локальный LLM, RAG (D3) или MCP. MCP подключать после стабилизации обычных Python-функций.
    \item SuiteSparse/Dziekonski и направление Regazzoni ROM вести как отдельные рабочие пакеты, не смешивать с текущим FEM-agent benchmark без отдельной гипотезы.
\end{enumerate}

\section{2026-11-10 - первый запуск с локальной Qwen}

\begin{enumerate}
    \item Пользователь установил Ollama и \texttt{qwen3.5:4b}; проверка \texttt{/api/tags} в его PowerShell подтвердила доступность модели. Изолированная среда проекта напрямую не видела локальный порт, поэтому запуски из проекта выполнялись с разрешённым пользователем доступом к \texttt{127.0.0.1:11434}. Системные настройки Windows не менялись.
    \item В Ollama-адаптер добавлены \texttt{think: false}, \texttt{num\_ctx=4096} и \texttt{num\_predict=128}, чтобы ограничить непрактично длинный ответ модели на короткие параметры задачи. Ruff и компиляционная проверка Python прошли.
    \item Реальный D4-запуск на запрос с \texttt{maxh=0.1}, quadratic elements корректно разобрал \texttt{poisson\_unit\_square\_smooth}, \texttt{maxh=0.1}, \texttt{order=2}; задача прошла. L2 error \texttt{1.7011e-4}, H1 seminorm error \texttt{1.2633e-2}, свободная невязка \texttt{2.10e-15}, 501 DoF; FEM solver занял 0.37 с. Один вызов модели занял 35.12 с. Отчёт: \texttt{results/agent\_runs/D4-20261110T145923993630Z/agent\_run.json}.
    \item Выполнено парное сравнение D1/D4 на одном фиксированном кейсе P1 <<default settings>>. D1 закончил расчёт и численную проверку успешно, но модель выбрала \texttt{maxh=0.1} вместо заданного в fixture значения \texttt{0.2}; потребовались три вызова модели, 151.08 с суммарно. Поэтому строгая проверка совпадения параметров отметила D1 как fail. D4 выбрал ожидаемые \texttt{maxh=0.2}, \texttt{order=2}, прошёл численную проверку одним вызовом модели за 35.64 с. Сравнение в целом помечено \texttt{all\_passed=false} из-за параметрического расхождения D1, а не из-за сбоя FEM. Отчёт: \texttt{results/agent\_benchmarks/agent-compare-20261110T150054102020Z/comparison.json}.
    \item Это пока один локальный прогон каждого варианта, не статистический вывод. Следующий шаг - сделать ожидание P1 явным в тексте запроса (убрать неоднозначное <<default settings>>), затем провести несколько парных повторов; отдельно ограничить число D1-вызовов/время и измерять корректность разбора, число вызовов, задержку, успешность верификации и численные метрики.
\end{enumerate}

\subsection{Повтор после уточнения P1}

\begin{enumerate}
    \item Формулировка P1 исправлена: вместо <<default settings>> текст теперь прямо задаёт \texttt{maxh=0.2} и quadratic elements. Повторное сравнение D1/D4 прошло полностью: оба выбрали ожидаемые параметры, прошли независимую численную проверку и получили одинаковые метрики (L2 \texttt{1.6335e-3}, свободная невязка \texttt{1.19e-15}, 125 DoF).
    \item D1 использовал 3 вызова модели за 158.99 с; D4 - 1 вызов за 34.85 с. Это только единичная пара на одном кейсе; преимущество по задержке пока предварительное.
    \item Отчёт: \texttt{results/agent\_benchmarks/agent-compare-20261110T150531948632Z/comparison.json}.
    \item Следующий шаг: запустить тот же явно заданный P1 несколько раз для оценки разброса, после этого расширить сравнение на P2-P5; отдельно проверить лимиты и недопустимые запросы, чтобы убедиться, что агент безопасно отклоняет неподдерживаемую физику.
\end{enumerate}

\section{2026-11-10 - P2-P5 и проверка границ поддержки}

\begin{enumerate}
    \item Запущено реальное сравнение D1/D4 на оставшихся кейсах P2-P5 с Ollama \texttt{qwen3.5:4b}. D4 точно воспроизвёл параметры и прошёл численный evaluator во всех четырёх кейсах. Начальная версия D1 прошла FEM и evaluator, но в P4 передала порядок 3 и maxh 0.1, при этом потеряла \texttt{grid\_samples=128} и \texttt{l2\_tolerance=0.03} из запроса; остальные P2, P3, P5 совпали с fixture. Полный исходный отчёт: \texttt{results/agent\_benchmarks/agent-compare-20261110T151057274661Z/comparison.json}.
    \item Исправлен D1 tool contract: \texttt{set\_task} теперь перечисляет все поля задачи, допустимые диапазоны и defaults; системная инструкция требует сохранять все явно запрошенные поля. P4 D1 повторён отдельно и прошёл со всеми четырьмя параметрами. L2 error \texttt{3.9001e-6}, свободная невязка \texttt{2.09e-15}, 1096 DoF. Отчёт: \texttt{results/agent\_benchmarks/agent-compare-20261110T153246027822Z/comparison.json}.
    \item По завершённым прогонам P1-P5, с исправленным контрактом D1, все пары D1/D4 точно воспроизвели параметры и прошли FEM evaluator. Но это один запуск на каждую пару, а не доказательство общей надёжности. D1 по-прежнему делает 3 вызова модели (около 149-212 секунд в этих примерах); D4 делает 1 (около 32-44 секунд).
    \item Выполнен контрольный запрос на 3D Navier-Stokes. До исправления D4 ошибочно объявил задачу \texttt{passed}, подменив её стандартной Poisson-задачей. Для обоих режимов добавлена детерминированная проверка области поддержки до LLM и FEM: запрос должен явно называть Poisson и не должен явно называть распространённые неподдерживаемые PDE (Navier-Stokes, Stokes, Maxwell, elasticity, Helmholtz, heat, wave). Иначе сохраняется \texttt{unsupported\_request}; LLM и решатель не вызываются.
    \item Добавлены проверки границ для обоих режимов и двух вариантов формулировки запроса. \texttt{pytest -q --basetemp .pytest-tmp tests/test\_agent\_cli\_scope.py tests/test\_workflow.py} завершился успешно: 6 passed; Ruff прошёл. Первичная ошибка pytest без \texttt{--basetemp} была вызвана запретом записи во временную папку Windows sandbox; повторный запуск с временным каталогом в проекте прошёл.
    \item Следующий шаг: повторить P1-P5 несколько раз на фиксированной версии модели, затем сравнить точность разбора, число обращений и задержку с разбросом. Лексическая проверка scope намеренно консервативна и пока покрывает только Poisson; при добавлении новых физических постановок её нужно заменить явным registry поддерживаемых task types.
    \item После обновления D1-контракта дополнительно прогнаны scope, D1 и D4 workflow unit tests: \texttt{7 passed}; Ruff и \texttt{compileall} прошли.
\end{enumerate}

\section{2026-11-10 - repeated P4 comparison}

\begin{enumerate}
    \item Добавлен параметр \texttt{--repeats} в \texttt{examples/evaluate\_agent\_architectures.py} (1-10): каждый прогон сохраняется отдельно; итоговый JSON содержит по каждой паре case/architecture pass rate, exact parameter-match rate, среднее число вызовов, среднее и медианное время. Команда описана в README.
    \item P4 прогнан трижды на D1 и трижды на D4 с Ollama \texttt{qwen3.5:4b}. Оба режима: 3/3 численно успешны и 3/3 раз точно восстановили параметры.
    \item D1: 3 вызова модели на запуск; среднее время 172.87 с, медиана 141.98 с (значения запусков 241.73, 141.98, 134.92 с). D4: 1 вызов; среднее 30.25 с, медиана 26.81 с (38.13, 26.81, 25.80 с). Наблюдаемый разброс указывает на необходимость большего числа повторов для устойчивой оценки задержки; все шесть текущих результатов успешны.
    \item Отчёт с отдельными прогонами и агрегатами: \texttt{results/agent\_benchmarks/agent-compare-20261110T154149619098Z/comparison.json}.
    \item Повторы ещё не выполнены на P1-P3 и P5. Следующий технический шаг - решить, нужна ли статистика для всех кейсов или P4 станет расширенным latency/reliability-кейсом, а P1-P5 останутся функциональным smoke benchmark; после решения можно экономно запускать дальнейшие повторы, не умножая без причины медленные D1 вызовы.
\end{enumerate}

\section{2026-11-10 - второй тип физической задачи: Stokes}

\begin{enumerate}
    \item Сверка с записью разговора с Александром: его исследовательская идея - найти/собрать agentic framework, который получает описание симуляции, формирует основной driver, запускает научное ПО и постобрабатывает результат; NGSolve/FEM - подходящая стабильная отправная точка. Он прямо предложил сначала изучить литературу и NGSolve, затем делать небольшие эксперименты с AI-агентами и NGSolve. Это exploratory research: он не требовал обучать новую LLM.
    \item Реализована первая ступень обобщения: allowlisted registry теперь маршрутизирует две типизированные задачи - manufactured Poisson и фиксированный Stokes tutorial вокруг цилиндра - к разным решателям и отдельным evaluators. Для Stokes проверяются конечность метрик, сходимость, L2 divergence, inlet/outlet mass balance и свободная невязка.
    \item \texttt{examples/stokes\_cylinder.py} переведён на ту же реализацию \texttt{src/ngsolve\_agent/stokes\_solver.py}, чтобы tutorial baseline и agent path не расходились. Реальный запуск при maxh=0.05: 3690 DoF, 760 элементов, divergence L2 \texttt{0.3156747}, inlet/outlet flux \texttt{0.41}, imbalance \texttt{0}, free residual \texttt{1.03e-14}.
    \item Реальный D4+Ollama запуск Stokes прошёл одним вызовом модели; модель выбрала ожидаемый Stokes task ID и maxh. Реальный D1+Ollama запуск после уплощения typed tool contract также прошёл, с тремя ожидаемыми решениями (set\_task/run\_solver/evaluate), без ошибки и восстановления. Первый D1-прогон выявил, что описания полей нельзя кодировать как псевдовложенные поля: модель попыталась передать ключ \texttt{Stokes fields}, схема отвергла его, затем D1 восстановился. Контракт исправлен, чистый повтор это подтвердил.
    \item После обобщения выполнен настоящий D4 Poisson regression при maxh=0.1/order=2: прошёл, L2 error \texttt{1.7011e-4}, residual \texttt{2.10e-15}.
    \item Полный тестовый набор: 18 passed; Ruff и compileall прошли. README описывает два поддерживаемых agent task types и текущие ограничения.
    \item Граница соответствия запросу Александра: уже есть реальный LLM-driven выбор задачи и параметров, вызов готового NGSolve driver, безопасный запуск фиксированного решателя, независимая оценка и постобработка/отчёт для двух задач. Пока агент НЕ генерирует новый основной driver-код по произвольному описанию; текущие Poisson/Stokes пути - закрытый prototype с allowlist. Следующий исследовательский этап - оценить D3: генерацию NGSolve driver с ограниченным набором шаблонов/инструментов, проверкой до исполнения и независимым evaluator; сравнить, какие запросы требуют новой генерации, а какие выразимы registry. До запуска сгенерированного кода нужен отдельный threat/sandbox design.
\end{enumerate}

\section{2026-11-10 - первый D3 driver-generation pilot и сравнение трёх архитектур}

\begin{enumerate}
    \item Добавлен D3. LLM генерирует Python-файл для выбора одного из двух поддерживаемых NGSolve task IDs и передачи параметров. До запуска AST-валидатор разрешает ровно два выражения: импорт \texttt{run} из фиксированного API и один вызов с простыми литералами. Затем исходник пересобирается из Pydantic-валидированных параметров, записывается как canonical driver и запускается отдельным Python \texttt{-I} процессом с минимальным набором переменных окружения. Именно исходный текст модели не исполняется.
    \item Первая D3 попытка на P4 была отклонена валидатором, потому что пример в prompt оказался двусмысленным и модель вложила параметры в \texttt{parameter=\{...\}}. Ошибку не обходили: уточнили prompt на плоские именованные аргументы, повторили и получили успешные P4 и S1 запуски. Оба создали candidate и canonical driver, оба прошли независимый evaluator.
    \item \texttt{examples/evaluate\_agent\_architectures.py} теперь поддерживает D1/D3/D4, общую P1-P5/S1 fixture \texttt{data/agent\_task\_cases.json}, повторяемые прогоны и сквозное время агента. End-to-end включает LLM, построение/проверку driver, subprocess, FEM и evaluator.
    \item Сводное сравнение D1/D3/D4 на P4 и S1 прошло полностью: 6/6 успешных результатов, 6/6 точных совпадений параметров. Сквозное время одного прогона:
    \begin{enumerate}
        \item P4: D1 143.6 s (3 LLM calls), D3 42.0 s (1 call), D4 51.9 s (1 call).
        \item S1: D1 128.0 s (3 calls), D3 50.4 s (1 call), D4 59.6 s (1 call).
    \end{enumerate}
    Один прогон на пару не даёт статистического вывода о превосходстве; D3 измерен только на этих двух кейсах.
    \item Сводный отчёт: \texttt{results/agent\_benchmarks/agent-compare-20261110T165408021049Z/comparison.json}. D3 candidate и canonical scripts лежат рядом с каждым запуском в подкаталогах \texttt{d3/}.
    \item Полный набор тестов после D3: 24 passed; Ruff, compileall и JSON fixture validation прошли.
\end{enumerate}

\subsection{Контрольная сверка с идеей Александра}

\begin{enumerate}
    \item Александр предложил исследовать agentic frameworks и проверить цепочку: описание симуляции -> основной driver -> запуск научного расчёта -> post-processing; начать с NGSolve/FEM и небольших экспериментов, не обязательно создавать или обучать свою LLM.
    \item Сейчас проверена архитектурная часть этой цепочки на двух задачах. D1 выбирает typed инструменты, D4 передаёт управление фиксированному workflow, D3 производит исполняемый main driver, но его содержание пока ограничено вызовом одного заранее реализованного solver API. Поэтому это proof-of-pipeline, а не генерация новой FEM-математики или произвольного решателя.
    \item Следующий этап: формализовать исследовательский вопрос и обзор framework/codegen подходов, затем расширять D3 от выбора готового solver к сборке driver из небольшой библиотеки проверенных FEM-компонентов. Сохранять независимые физические checks, ограниченный AST, воспроизводимые запуски и сравнение D1/D3/D4. Сначала - paraphrase/robustness и invalid request tests; потом новые solver components. Не заявлять научную новизну до сравнения с опубликованными agentic scientific-computing системами.
\end{enumerate}

\section{2026-11-10 - сверка с гипотезами внешнего ревью по D3}

\begin{enumerate}
    \item Получен текст ревью, в котором автор прямо указал, что исходники и comparison report ему не передали; его security/benchmark замечания были гипотезами, не проверкой кода.
    \item По фактической реализации подтверждено: Python \texttt{-I} не является OS sandbox; однако D3 не исполняет исходную строку модели. AST допускает только точный import фиксированного \texttt{run} API, один вызов и простые литеральные keyword arguments; параметры валидируются Pydantic, из них строится новый canonical driver, который и запускается. Подпроцесс получает небольшой env allowlist без \texttt{NGSOLVE\_LLM\_API\_KEY}; произвольные файлы/сеть по OS не ограничены, но исходник модели не имеет выражений, которыми мог бы ими воспользоваться. Уязвимость в самом доверенном solver/dependency этим не закрывается.
    \item Подтверждён ресурсный пробел: \texttt{maxh > 0} допускал патологически мелкую сетку. Заданы нижние пределы \texttt{maxh >= 0.025} для Poisson и Stokes; ограничения описаны в prompt и схемах.
    \item AST парсер уже имел лимит исходника 16,000 знаков, точную форму из двух statements и отказ от неразрешённых imports/calls/kwargs. Добавлены обработка parser recursion/memory errors и лимит 128 AST nodes. D3 subprocess уже имел timeout 600 секунд.
    \item Подтверждён статистический вывод ревью: P4/S1 D1/D3/D4 имеют по одному запуску в общем сравнении; они не дают надёжной оценки частоты отказов. Данных недостаточно для claims о превосходстве архитектуры. D1/D4 ранее повторялись три раза только на P4; D3 - один раз.
    \item Полный тестовый набор после resource/AST изменений: 27 passed; Ruff прошёл. Следующая статистическая работа - сначала несколько повторов D3 на обеих задачах и paraphrase cases; большой 20-run sweep отложен из-за длительных локальных D1 прогонов.
    \item Так как Клод не видел файлы, собран узкий zip для повторного независимого ревью: содержит необходимые модули, тесты, P4/S1 comparison report, candidate/canonical drivers и точное задание ревью. Путь: \texttt{results/review\_bundles/claude\_d3\_review\_20261110.zip}.
\end{enumerate}

\section{2026-11-10 - проверка ответа Клода и paraphrase robustness}

\begin{enumerate}
    \item Исходный ответ Клода пришёл без приложенных исходников, поэтому он сам обозначил все пункты как гипотезы. По коду проверено: D3 не исполняет raw model source; точная AST-форма и typed параметры обязательны, canonical driver пересобирается; subprocess получает env allowlist без LLM API key. \texttt{-I} не OS sandbox, поэтому код доверенных solver/dependencies всё ещё имеет права текущего пользователя; D3 не даёт произвольному тексту модели доступ к командам/файлам.
    \item Закрыт подтверждённый риск слишком малого \texttt{maxh}: схемы ограничивают Poisson и Stokes интервалом от 0.025; нижняя граница синхронизирована с LLM prompts. Добавлены parser recursion/memory handling и лимит 128 AST nodes (source limit 16 KB уже был).
    \item Входной scope gate распознал единицу-квадрат только по словам. Новая paraphrase с \texttt{[0,1]\^{}2} до исправления была отклонена до LLM; добавлено распознавание стандартной записи. После этого D3 на P6 восстановил все параметры и прошёл solve/evaluator. Stokes paraphrase S2 также прошёл. Это зафиксировано в \texttt{results/agent\_benchmarks/agent-compare-20261110T171531076845Z/} и \texttt{results/agent\_benchmarks/agent-compare-20261110T171932139447Z/}.
    \item Полный suite после обновлений: 28 passed; Ruff, compileall и JSON fixture validation passed.
    \item Обновлён zip для Клода: \texttt{results/review\_bundles/claude\_d3\_review\_20261110.zip}. В нём текущие исходники/тесты, точный review request, основные сравнения, paraphrase reports и generated candidate/canonical drivers. Клоду нужно загрузить этот архив; предыдущий ответ был без файлов.
    \item Статистическая осторожность сохраняется: у D3 пока нет повторов на фиксированном кейсе P4, поэтому скорость/надёжность D3 против D1/D4 не объявлять доказанными.
\end{enumerate}

\section{2026-11-11 - независимый запуск Клода и фиксация evaluator policy}

\begin{enumerate}
    \item Получено новое ревью с архивом исходников и независимым запуском solver/evaluator в Linux. Клод сообщил о 30,000 мутациях AST-кандидатов без обхода текущего узкого allowlist; его численные результаты Poisson P4 и Stokes на нескольких сетках полезны как наблюдения, но ревью не запускало полный LLM agent CLI и не проверяло Windows subprocess.
    \item Подтверждён критичный дефект: модель могла передать \texttt{l2\_tolerance}, \texttt{divergence\_tolerance}, \texttt{flux\_imbalance\_tolerance} и \texttt{residual\_tolerance}, тем самым менять критерий собственного успеха. Эти поля удалены из Pydantic task schemas, D1/D3/D4 prompts/contracts и benchmark fixtures. P4/P6 больше не просят настраивать порог точности.
    \item Evaluator использует фиксированные пороги: Poisson L2 error \texttt{<= 2e-2}, Stokes relative flux imbalance \texttt{<= 1e-8}, free-DOF residual \texttt{<= 1e-10}. Для Stokes divergence остаётся метрикой, но больше не служит pass/fail gate: для tutorial-задачи нет заданного здесь эталонного решения, а прежний порог был произвольным. Документация отмечает, что эти критерии - pilot checks, а не доказательство точности.
    \item Добавлены проверки контракта: удалённые evaluator thresholds отвергаются схемами; высокий Stokes divergence сам по себе не валит задачу; превышение фиксированного Poisson порога валит её. Тесты и численные baseline после этого изменения ещё требуется прогнать в локальном Windows окружении; текущая запись не утверждает, что этот запуск уже состоялся.
    \item Следующий технический этап: прогнать тесты и agent cases после смены evaluator; затем сделать независимую Poisson сеточную/order convergence серию и подобрать Stokes reference/manufactured контроль до расширения benchmark и заявлений о научном качестве.
\end{enumerate}

\section{2026-11-11 - полный пакет независимого ревью}

\begin{enumerate}
    \item Выяснена причина замечания Клода о неполном архиве: предыдущий zip был собран как узкий D3 security review bundle и не включал package entry points, lock-файл, все примеры, полный development log и прогоны D1/D4.
    \item Собран полный пакет \texttt{results/review\_bundles/claude\_review\_complete\_20261111.zip}: весь \texttt{src/}, \texttt{examples/}, \texttt{tests/}, \texttt{data/}, README, \texttt{DEVELOPMENT\_LOG.md}, \texttt{pyproject.toml}, \texttt{uv.lock}, а также JSON/driver artifacts P4 3-run D1/D4, paired P4/S1 D1/D3/D4, P6 и S2 paraphrases.
    \item В пакет добавлен manifest и обновлённый запрос ревью. В них явно сказано, что historical reports предшествуют новым фиксированным evaluator-порогам и не являются результатом повторного запуска текущей версии.
    \item Следующие приоритеты из отчёта: (1) запустить текущую версию и перепроверить baseline; (2) добавить регулярный baseline и отрицательные/неоднозначные запросы; (3) прогнать P6/S2 через D1/D4; (4) сделать D3 convergence-study case, где driver действительно задаёт серию расчётов и post-processing наблюдаемого порядка; (5) закрыть subprocess tests/output limits и LLM token cap.
\end{enumerate}

\end{document}